\subsection{Eigenvalues, Eigenspaces, and Spectral Criteria}

Eigenstructure begins by identifying the scalars for which a shifted operator has a nontrivial kernel.

Throughout this subsection, $\mathcal{V}$ is a vector space over $\mathbb{F}$, and $n\in\mathbb{Z}_{>0}$ whenever an $n\times n$ matrix occurs.

\begin{definition} \label{def:eigenvalue-eigenvector} \hlt{Eigenvalue and Eigenvector}\\
Let $T\in\LL(\mathcal{V})$.
\begin{enumerate}[label=(\roman*)]
\setlength{\itemsep}{0pt}
\item A scalar $\lambda\in\mathbb{F}$ is an \hlt{eigenvalue} of $T$ iff $(\exists\,\mathbf{v}\in\mathcal{V}\setminus\{\mathbf{0}\})\,T\mathbf{v}=\lambda\mathbf{v}$.
\item For $\lambda\in\mathbb{F}$, a vector $\mathbf{v}\in\mathcal{V}$ is an \hlt{eigenvector} of $T$ corresponding to $\lambda$ iff $\mathbf{v}\neq\mathbf{0}\land T\mathbf{v}=\lambda\mathbf{v}$.
\end{enumerate}
\end{definition}

\begin{definition} \label{def:eigenspace} \hlt{Eigenspace}\\
Let $T\in\LL(\mathcal{V})$ and $\lambda\in\mathbb{F}$. The \hlt{eigenspace} of $T$ corresponding to $\lambda$ is $E_{\lambda}(T)\defeq\ker(T-\lambda I_{\mathcal{V}})=\{\mathbf{v}\in\mathcal{V}:T\mathbf{v}=\lambda\mathbf{v}\}$.\\
Equivalently, $\lambda$ is an eigenvalue of $T$ iff $E_{\lambda}(T)\neq\{\mathbf{0}\}$.
\end{definition}

\begin{definition} \label{def:spectrum-and-geometric-multiplicity} \hlt{Spectrum and Geometric Multiplicity}\\
Let $T\in\LL(\mathcal{V})$.
\begin{enumerate}[label=(\roman*)]
\setlength{\itemsep}{0pt}
\item The \hlt{spectrum} of $T$ is $\operatorname{spec}(T)\defeq\{\lambda\in\mathbb{F}:E_{\lambda}(T)\neq\{\mathbf{0}\}\}$.
\item If $\mathcal{V}$ is finite-dimensional and $\lambda\in\operatorname{spec}(T)$, then the \hlt{geometric multiplicity} of $\lambda$ is $\dim E_{\lambda}(T)$.
\item For $A\in\operatorname{Mat}(n,n,\mathbb{F})$ and $\lambda\in\mathbb{F}$, set $E_{\lambda}(A)\defeq\ker(A-\lambda I_n)$. Define $\operatorname{spec}(A)\defeq\{\mu\in\mathbb{F}:E_{\mu}(A)\neq\{\mathbf{0}\}\}$.
\end{enumerate}
\end{definition}

\begin{theorem} \label{thm:equivalent-eigenvalue-conditions} \hlt{Equivalent Conditions for an Eigenvalue}\\
Let $\mathcal{V}$ be finite-dimensional, let $T\in\LL(\mathcal{V})$, and let $\lambda\in\mathbb{F}$. The following are equivalent:
\begin{enumerate}[label=(\roman*)]
\setlength{\itemsep}{0pt}
\item $\lambda$ is an eigenvalue of $T$;
\item $T-\lambda I_{\mathcal{V}}$ is not injective;
\item $T-\lambda I_{\mathcal{V}}$ is not surjective;
\item $T-\lambda I_{\mathcal{V}}$ is not invertible.
\end{enumerate}
\begin{proof}
Set $S\defeq T-\lambda I_{\mathcal{V}}$.
\begin{enumerate}[label=(\roman*)]
\setlength{\itemsep}{0pt}
\item[$(i)$] $\Rightarrow (ii)$ Choose $\mathbf{v}\in\mathcal{V}\setminus\{\mathbf{0}\}$ with $T\mathbf{v}=\lambda\mathbf{v}$ by (i).\\
The membership $\mathbf{v}\in\ker(S)\setminus\{\mathbf{0}\}$ holds by $S\mathbf{v}=T\mathbf{v}-\lambda\mathbf{v}=\mathbf{0}$, and $S$ is not injective by Proposition \ref{prop:injective-iff-kernel-zero}.
\item[$(ii)$] $\Rightarrow (iii)$ The operator $S$ is not surjective by Theorem \ref{thm:finite-dimensional-operator-invertible-equivalences}(ii),(iii).
\item[$(iii)$] $\Rightarrow (iv)$ The operator $S$ is not invertible by Theorem \ref{thm:finite-dimensional-operator-invertible-equivalences}(iii),(i).
\item[$(iv)$] $\Rightarrow (i)$ The operator $S$ is not injective by Theorem \ref{thm:finite-dimensional-operator-invertible-equivalences}(i),(ii), and $\ker(S)\neq\{\mathbf{0}\}$ by Proposition \ref{prop:injective-iff-kernel-zero}.\\
Choose $\mathbf{v}\in\ker(S)\setminus\{\mathbf{0}\}$. The equality $T\mathbf{v}=\lambda\mathbf{v}$ holds by $S\mathbf{v}=\mathbf{0}$. The scalar $\lambda$ is an eigenvalue of $T$ by the choice of $\mathbf{v}$.
\end{enumerate}
\hlt{Structural Commentary}: The kernel test links singular shifts to the finite-dimensional invertibility cycle.\\
\hlt{Key Insight}: A scalar is spectral exactly when its shifted operator has nonzero kernel.
\end{proof}
\end{theorem}

\begin{theorem} \label{thm:distinct-eigenvalues-independent-eigenvectors} \hlt{Eigenvectors for Distinct Eigenvalues Are Independent}\\
Let $m\in\mathbb{Z}_{>0}$ and $T\in\LL(\mathcal{V})$. Suppose $\lambda_1,\ldots,\lambda_m$ are pairwise distinct eigenvalues of $T$ with corresponding eigenvectors $\mathbf{v}_1,\ldots,\mathbf{v}_m$.\\
Then $(\mathbf{v}_1,\ldots,\mathbf{v}_m)$ is linearly independent.
\begin{proof}
For $m=1$, the list is independent by $\mathbf{v}_1\neq\mathbf{0}$. Assume $m\geq2$ and suppose $(\mathbf{v}_1,\ldots,\mathbf{v}_m)$ is dependent.\\
Let $k$ be least index for which $(\mathbf{v}_1,\ldots,\mathbf{v}_k)$ is dependent. Bounds $k\leq m$ and $k\geq2$ hold by choice of $k$ and $\mathbf{v}_1\neq\mathbf{0}$, respectively.\\
The list $(\mathbf{v}_1,\ldots,\mathbf{v}_{k-1})$ is independent by minimality of $k$. Choose $c_1,\ldots,c_k\in\mathbb{F}$, not all zero, with $\sum\limits_{i=1}^{k}c_i\mathbf{v}_i=\mathbf{0}$ by dependence.\\
The coefficient $c_k\neq0$ by independence of the preceding prefix: $c_k=0$ would force $c_1=\cdots=c_{k-1}=0$.\\
Set $a_i\defeq-\frac{c_i}{c_k}$ $(1\leq i<k)$. The representation $\mathbf{v}_k=\sum\limits_{i=1}^{k-1}a_i\mathbf{v}_i$ holds by the chosen dependence relation.\\
The equality $\mathbf{0}=(T-\lambda_kI_{\mathcal{V}})\mathbf{v}_k=\sum\limits_{i=1}^{k-1}a_i(\lambda_i-\lambda_k)\mathbf{v}_i$ holds by linearity and the eigenvector identities.\\
The relation $(\forall\,i\in\{1,\ldots,k-1\})\,a_i(\lambda_i-\lambda_k)=0$ holds by independence of the preceding prefix.\\
The relation $(\forall\,i\in\{1,\ldots,k-1\})\,a_i=0$ holds by pairwise distinctness of the eigenvalues.\\
The equality $\mathbf{v}_k=\mathbf{0}$ holds by its representation and the preceding coefficient identities, contrary to the eigenvector hypothesis.\\
The original list is linearly independent by contradiction.\\
\hlt{Structural Commentary}: Minimal dependence exposes the last coefficient; the spectral shift isolates the rest.\\
\hlt{Key Insight}: The shift by $\lambda_k$ rescales every earlier eigenvector nontrivially.
\end{proof}
\end{theorem}

\begin{corollary} \label{cor:number-of-distinct-eigenvalues} \hlt{Number of Distinct Eigenvalues}\\
If $\mathcal{V}$ is finite-dimensional and $T\in\LL(\mathcal{V})$, then $\abs{\operatorname{spec}(T)}\leq\dim(\mathcal{V})$.
\begin{proof}
Set $d\defeq\dim(\mathcal{V})$, and choose a basis $\beta$ of length $d$ by finite dimensionality and the definition of $d$.\\
Suppose $\abs{\operatorname{spec}(T)}>d$. Choose distinct $\lambda_1,\ldots,\lambda_{d+1}\in\operatorname{spec}(T)$.\\
Choose corresponding eigenvectors $\mathbf{v}_1,\ldots,\mathbf{v}_{d+1}$ by the definition of the spectrum.\\
The list $(\mathbf{v}_1,\ldots,\mathbf{v}_{d+1})$ is independent by Theorem \ref{thm:distinct-eigenvalues-independent-eigenvectors}, while $\beta$ spans $\mathcal{V}$.\\
The inequality $d+1\leq d$ holds by Theorem \ref{thm:finitedimlenofvec}, a contradiction. The claimed bound holds by contradiction.\\
\hlt{Structural Commentary}: Spectral representatives form an independent list bounded in length by a basis.\\
\hlt{Key Insight}: Each distinct eigenvalue requires a separate independent direction.
\end{proof}
\end{corollary}

\begin{corollary} \label{cor:invertible-iff-zero-not-eigenvalue} \hlt{Invertibility and the Zero Eigenvalue}\\
Let $\mathcal{V}$ be finite-dimensional and let $T\in\LL(\mathcal{V})$. Then $T$ is invertible iff $0\notin\operatorname{spec}(T)$.
\begin{proof}
\begin{enumerate}[label=\roman*.]
\setlength{\itemsep}{0pt}
\item[$(\Rightarrow)$] Suppose $T$ is invertible. The scalar $0$ is not an eigenvalue of $T$ by Theorem \ref{thm:equivalent-eigenvalue-conditions}(iv),(i).\\
The exclusion $0\notin\operatorname{spec}(T)$ holds by the definition of $\operatorname{spec}(T)$.
\item[$(\Leftarrow)$] Suppose $0\notin\operatorname{spec}(T)$. The scalar $0$ is not an eigenvalue of $T$ by the definition of $\operatorname{spec}(T)$.\\
The operator $T=T-0I_{\mathcal{V}}$ is invertible by Theorem \ref{thm:equivalent-eigenvalue-conditions}(i),(iv).
\end{enumerate}
\hlt{Structural Commentary}: At zero, the singular-shift criterion equates invertibility with spectral exclusion.\\
\hlt{Key Insight}: The zero eigenspace is $\ker(T)$.
\end{proof}
\end{corollary}

\begin{corollary} \label{cor:matrix-invertible-iff-zero-not-spectrum} \hlt{Matrix Invertibility and the Spectrum}\\
Let $A\in\operatorname{Mat}(n,n,\mathbb{F})$. Then $A$ is invertible iff $0\notin\operatorname{spec}(A)$.
\begin{proof}
\begin{enumerate}[label=\roman*.]
\setlength{\itemsep}{0pt}
\item[$(\Rightarrow)$] Suppose $A$ is invertible. The equality $\operatorname{rank}(A)=n$ holds by Corollary \ref{cor:square-matrix-invertible-iff-full-rank}.\\
The equality $\dim(\ker(A))=0$ holds by Corollary \ref{cor:matrix-rank-nullity}(i).\\
The relations $\{\mathbf{0}\}\leq\ker(A)$ and $\dim(\{\mathbf{0}\})=\dim(\ker(A))=0$ hold by $A\mathbf{0}=\mathbf{0}$ and the preceding equality.\\
The identity $\ker(A)=\{\mathbf{0}\}$ holds by Proposition \ref{prop:dimension-bounds-and-basis-criteria}(ii). The exclusion $0\notin\operatorname{spec}(A)$ holds by the matrix-spectrum definition.
\item[$(\Leftarrow)$] Suppose $0\notin\operatorname{spec}(A)$. The identity $\ker(A)=\{\mathbf{0}\}$ holds by the matrix-spectrum definition.\\
Equality $\dim(\ker(A))=0$ holds by $\ker(A)=\{\mathbf{0}\}$, $\operatorname{rank}(A)=n$ by Corollary \ref{cor:matrix-rank-nullity}(i). Matrix $A$ is invertible by Corollary \ref{cor:square-matrix-invertible-iff-full-rank}.
\end{enumerate}
\hlt{Structural Commentary}: Rank--nullity carries full rank to the zero-kernel matrix-spectrum test.\\
\hlt{Key Insight}: Zero is spectral exactly when $\ker(A)\neq\{\mathbf{0}\}$.
\end{proof}
\end{corollary}
